\section{Experimental Setup}
\label{sec:experimental-setup}

\subsection{A note on the dataset behind this report}
\label{sec:dataset-note}

The dataset this section and Section~\ref{sec:results} analyze, committed at \texttt{results/sample\_run/}, is a real measurement: the custom ring and the vendor \texttt{MPI\_Allreduce} were run on the workstation described below, under a working Microsoft MPI installation, and the resulting per-configuration timings are what every figure and table here is derived from. Before those timings were trusted, correctness was established independently: the full \texttt{ctest} matrix (Section~\ref{sec:implementation}), including non-power-of-two process counts and the zero-size-chunk cases, passes against that same real MPI, and each configuration's ring output is checked against both vendor \texttt{MPI\_Allreduce} and an independently computed reference before any timing of it is recorded.

Two honest scoping caveats follow from the environment and belong here rather than in a footnote. First, this is a single-node measurement: all ranks share one machine, so the transport underneath the point-to-point calls is Microsoft MPI's shared-memory path, not a network fabric, and the fitted latency and bandwidth constants of Section~\ref{sec:cost-model} describe that path specifically. Second, Microsoft MPI exposes no equivalent of Open~MPI's \texttt{coll/tuned} algorithm-selection controls, so the forced-ring \texttt{MPI\_Allreduce} variant that an apples-to-apples ring-versus-ring comparison would use (Section~\ref{sec:mca-forcing}) cannot be produced here; this report therefore compares the custom ring against the vendor's default \texttt{MPI\_Allreduce} selection only. Both gaps are addressed by the same future step, a multi-node sweep under Open~MPI via \texttt{scripts/run\_slurm\_sweep.sbatch}, which the analysis pipeline consumes without modification.

\subsection{Target hardware and software}

The benchmarking machine is a workstation with an Intel Core i7-14700K (8 performance cores plus 12 efficiency cores, 28 hardware threads, and a 33~MB L3 cache that turns out to matter for the results below), 32~GB of DDR5 memory, and, unused by this CPU-only, MPI-only project, an NVIDIA RTX~5070. The measurements here were collected under the Microsoft MPI 10.1 runtime, with the MSYS2 UCRT64 \texttt{mpicxx} wrapper and GCC~15.2 providing the MPI-2 development headers and link library, built out-of-source with CMake under \texttt{-Wall -Wextra -Wpedantic} treated as errors. The build specification also names Open~MPI \citep{gabriel2004openmpi} 4.x or 5.x as the intended cluster runtime; that route is what enables the forced-ring comparison of Section~\ref{sec:mca-forcing}. Running \texttt{check\_mpi\_\allowbreak algorithms.sh} in \texttt{scripts/} on an Open~MPI machine prints its version and the numeric algorithm IDs \texttt{coll\_tuned\_allreduce\_algorithm} enumerates (which have shifted across Open~MPI releases), and \texttt{docs/BENCHMARKING.md} is where those values get recorded next to whatever sweep they came from.

\subsection{Sweep ranges and methodology}

The message-size sweep covers powers of two in bytes from $8=2^3$ through $134217728=2^{27}$ (128~MiB in the power-of-two sense \texttt{nccl-tests} uses \citep{nvidia2024ncclperf}), 25 sizes. The process-count sweep for the run analyzed here uses $N \in \{2,4,8,16\}$, the harness's \texttt{--quick} grid, chosen to keep total wall-clock time reasonable on a shared single-node desktop while still spanning small, moderate, and (relative to this machine's core count) contended process counts; the full $2..16$ range, including the non-power-of-two counts that Section~\ref{sec:step-count}'s step-count argument applies to identically, is available from the same harness with the flag omitted. Because $N$ is fixed for the lifetime of an MPI job, the sweep launcher (\texttt{scripts/run\_local\_sweep.ps1} on this Windows desktop, \texttt{scripts/run\_local\_sweep.sh} on Unix) starts one \texttt{mpiexec} per $N$ value; \texttt{apps/benchmark} itself sweeps every message size and every requested algorithm inside that one launch, so MPI startup and teardown overhead is paid once per $N$, not once per (size, algorithm) combination.

Each configuration is preceded by an \texttt{MPI\_Barrier} and a fixed number of untimed warm-up iterations (so first-touch page faults and any connection setup are not measured), then times a variable number of repetitions capped by a wall-clock budget rather than a fixed count, so small, fast configurations can accumulate many repetitions while the largest, slowest ones do not make the sweep take unreasonably long. Each repetition's elapsed time is measured with \texttt{MPI\_Wtime} on every rank and reduced to rank~0 with \texttt{MPI\_Reduce(..., MPI\_MAX, ...)}, since an allreduce is only as fast as its slowest participating rank, not as fast as whichever rank happens to be timing it. Min, mean, median, max, and standard deviation are all recorded per configuration; this report and its analysis pipeline use the \textbf{median} as "achieved" performance throughout, since it is far less sensitive than the minimum to the occasional favorable scheduling accident, while remaining far less sensitive than the mean to the occasional unfavorable one.

\subsection{Isolating implementation quality from algorithm choice}
\label{sec:mca-forcing}

The comparison this report can draw on a Microsoft MPI desktop is the custom ring against the vendor's default \texttt{MPI\_Allreduce} selection, which is the baseline an application actually gets. There is a second, sharper comparison that this environment cannot produce but an Open~MPI cluster can, and it is worth stating precisely because Section~\ref{sec:discussion} interprets the default-selection result partly by reference to it. Under Open~MPI, \texttt{MPI\_Allreduce} can be forced to the vendor's own ring algorithm via \texttt{--mca coll\_tuned\_use\_\allowbreak dynamic\_rules 1 \allowbreak --mca coll\_tuned\_\allowbreak allreduce\_algorithm <id>} on the \texttt{mpirun} invocation itself, giving a ring-versus-ring measurement that isolates this implementation's quality from the vendor's algorithm-selection logic: any remaining gap is implementation, not algorithm. The numeric algorithm id is read from \texttt{check\_mpi\_\allowbreak algorithms.sh}'s output on the target machine rather than assumed, since it has moved across Open~MPI releases, and whatever id and Open~MPI version that sweep runs against belongs in \texttt{docs/BENCHMARKING.md} next to its results.

Because MCA parameters are process-environment state read once at \texttt{MPI\_Init}, forcing the ring algorithm is a property of how \texttt{mpirun} was invoked, not something \texttt{apps/benchmark} can change mid-run; on Open~MPI the sweep script therefore performs two launches per $N$, one plain and one with the forcing flags, appending both to the same results file. Microsoft MPI has no counterpart to these controls, so the run analyzed here emits the ring and default-selection rows only, and the forced-ring column is left for the cluster sweep.
